\chapter{稳恒电流}

在静电学中，我们研究电荷如何建立电场与电势；而在电路问题中，我们更关心电荷如何\emph{持续流动}。一旦导体两端存在稳定的电势差，自由电荷便会在导体内部形成定向移动，从而产生电流。宏观上看，电流是单位时间内通过某一截面的电荷量；微观上看，电流则来自大量载流子在热运动背景上的微弱漂移。由此出发，我们可以自然地建立欧姆定律、基尔霍夫定律，并进一步研究带有电容、电感元件的暂态过程。

本章的主线有两条：
\begin{itemize}
  \item 从\emph{连续变化}的角度理解电路量：$I=\dv*{Q}{t}$，功率 $P=\dv*{W}{t}$；
  \item 把高中常见电路结论看成\emph{微分方程的解}：RC、RL 电路的指数规律都来自一阶线性方程。
\end{itemize}

\section{电流与电流密度}

\subsection{电流的定义}

\begin{definition}
设在时间区间 $[t,t+\dd t]$ 内，有电量 $\dd Q$ 通过导体某一横截面，则电流定义为
\[
I=\dv{Q}{t}.
\]
若在有限时间 $\Delta t$ 内通过的电量为 $\Delta Q$，则平均电流为
\[
I_{\text{avg}}=\frac{\Delta Q}{\Delta t}.
\]
\end{definition}

这个定义说明：电流并不是“某处储存了多少电”，而是“电荷通过截面的快慢”。因此，若已知累积通过截面的电量函数 $Q(t)$，只需对时间求导便可得到瞬时电流；反过来，若已知电流函数，则有
\[
Q(t)-Q(t_0)=\int_{t_0}^{t} I(\tau)\,\dd \tau.
\]

\begin{remark}
在金属导体中，实际运动的是自由电子，它们带负电，运动方向与约定的电流方向相反。电路分析中默认采用\emph{正电荷运动方向}作为电流正方向。
\end{remark}

\begin{example}
某导线中通过某截面的电量满足
\[
Q(t)=2t^3-3t^2+6t \qquad (\text{单位：C}).
\]
求 $t=2\,\text{s}$ 时的瞬时电流，以及从 $t=0$ 到 $t=2\,\text{s}$ 的平均电流。

\textbf{解：}
由定义
\[
I(t)=\dv{Q}{t}=6t^2-6t+6.
\]
故
\[
I(2)=24-12+6=18\,\text{A}.
\]
平均电流为
\[
I_{\text{avg}}=\frac{Q(2)-Q(0)}{2-0}=\frac{16-0}{2}=8\,\text{A}.
\]
可见平均电流与瞬时电流一般并不相同。
\end{example}

\subsection{电流密度}

如果只知道总电流，还无法描述电荷在导体内部是如何分布流动的。为此需要引入矢量量——电流密度。

\begin{definition}
设载流子数密度为 $n$，每个载流子电荷量为 $q$，其定向漂移速度为 $\vb{v}_d$，则电流密度定义为
\[
\vb{J}=nq\vb{v}_d.
\]
若导体截面的单位法向量为 $\uvec{n}$，则通过面积元 $\dd S$ 的电流元为
\[
\dd I=\vb{J}\cdot \uvec{n}\,\dd S.
\]
从而总电流为
\[
I=\int_S \vb{J}\cdot \dd \vb{S}.
\]
\end{definition}

若导体为截面积均匀的直导线，且 $\vb{J}$ 在截面上均匀分布并平行于导线轴线，则上式退化为
\[
I=JS,
\]
其中 $S$ 为横截面积，$J=\magn{\vb{J}}$ 为电流密度大小。

\begin{remark}
电流密度是局域量，类似于力学中的质量密度。总电流是对整个截面的积分结果，而电流密度描述的是“每单位面积上有多强的电流穿过”。
\end{remark}

\begin{example}
一根半径为 $r$ 的圆柱形导线中，电流密度沿半径分布为
\[
J(\rho)=J_0\left(1-\frac{\rho^2}{r^2}\right), \qquad 0\le \rho\le r.
\]
求总电流 $I$。

\textbf{解：}
取极坐标面积元 $\dd S=2\pi \rho\,\dd \rho$，则
\[
I=\int_S J\,\dd S
=\int_0^r J_0\left(1-\frac{\rho^2}{r^2}\right)2\pi\rho\,\dd \rho.
\]
计算得
\[
I=2\pi J_0\int_0^r \left(\rho-\frac{\rho^3}{r^2}\right)\dd \rho
=2\pi J_0\left(\frac{r^2}{2}-\frac{r^2}{4}\right)
=\frac{\pi J_0 r^2}{2}.
\]
因此，虽然中心处电流密度最大，但总电流仍然由对整个截面的积分决定。
\end{example}

\subsection{稳恒电流与电荷守恒}

所谓\emph{稳恒电流}，是指导体中各处电流不随时间改变。在稳恒条件下，导体内部电荷不会无限积累，否则局部电场将随时间变化，电流也就不再稳定。

电荷守恒的基本思想是：对任何一个封闭体积，内部电荷的减少率等于流出边界的电流通量。这一思想在电路节点处表现为“流入电流之和等于流出电流之和”，稍后将发展为基尔霍夫电流定律。

\begin{law}[电荷守恒的电路表述]
对任一节点，在任意时刻都有
\[
\sum I_{\text{入}}=\sum I_{\text{出}}.
\]
若规定流入为正、流出为负，则可写为
\[
\sum_k I_k=0.
\]
\end{law}

\section{欧姆定律的微观解释}

\subsection{漂移速度}

金属中的自由电子始终在做无规则热运动，但这种热运动的平均速度为零，不产生宏观电流。当外电场存在时，电子在碰撞间隔内受到电场力作用，会叠加一个很小的定向平均速度，这就是\emph{漂移速度}。

设载流子质量为 $m$，电荷量为 $q$，平均碰撞时间为 $\tau$。在外电场 $\vb{E}$ 中，载流子加速度近似为
\[
\vb{a}=\frac{q\vb{E}}{m}.
\]
在一次碰撞前的平均自由运动时间约为 $\tau$，于是漂移速度满足近似关系
\[
\vb{v}_d\approx \frac{q\tau}{m}\vb{E}.
\]
代入电流密度公式即得
\[
\vb{J}=nq\vb{v}_d=nq\left(\frac{q\tau}{m}\vb{E}\right)
=\frac{nq^2\tau}{m}\vb{E}.
\]

\begin{definition}
若材料满足
\[
\vb{J}=\sigma \vb{E},
\]
则称 $\sigma$ 为该材料的\emph{电导率}。其倒数
\[
\rho=\frac{1}{\sigma}
\]
称为\emph{电阻率}，于是也可写成
\[
\vb{E}=\rho\vb{J}.
\]
\end{definition}

上式就是欧姆定律的微观形式。它表明：欧姆导体中，电流密度与电场强度成正比，比例系数由材料本身决定。

\begin{law}[欧姆定律]
对一段温度近似不变的均匀导体，电压与电流满足
\[
U=IR.
\]
若导体长度为 $\ell$、截面积为 $S$，则其电阻为
\[
R=\rho\frac{\ell}{S}.
\]
\end{law}

\subsection{从微观形式到宏观形式}

对均匀直导体，内部电场近似沿轴线方向且大小处处相同，于是
\[
U=E\ell,
\qquad
I=JS.
\]
由微观欧姆定律 $E=\rho J$ 得
\[
U=\rho J\ell=\rho\frac{I}{S}\ell,
\]
从而
\[
U=I\left(\rho\frac{\ell}{S}\right).
\]
因此宏观欧姆定律不过是微观关系在一段均匀导体上的积分结果。

作为一个直接应用，一根铜导线长 $20\,\text{m}$，截面积 $1.0\times 10^{-6}\,\text{m}^2$，若铜的电阻率取 $1.7\times 10^{-8}\,\Omega\cdot\text{m}$，则
\[
R=\rho\frac{\ell}{S}=1.7\times 10^{-8}\times\frac{20}{1.0\times 10^{-6}}=0.34\,\Omega.
\]
当电流为 $2\,\text{A}$ 时，导线两端电压为 $U=IR=0.68\,\text{V}$，故平均电场强度
\[
E=\frac{U}{\ell}=\frac{0.68}{20}=3.4\times 10^{-2}\,\text{V/m}.
\]
这说明维持导体中稳恒电流所需的电场往往并不强。

\begin{remark}
漂移速度通常远小于电子热运动速度，也远小于电信号传播速度。电路闭合后灯泡几乎立刻发光，并不是因为电子从电源“跑到”灯泡很快，而是因为电场在整个导体回路中迅速建立起来。
\end{remark}

\section{基尔霍夫定律}

复杂电路中，仅靠单个元件的欧姆定律通常不足以确定各支路电流和各点电势，还需要借助电荷守恒与能量守恒建立整体方程。

\subsection{节点定律}

\begin{law}[基尔霍夫第一定律（节点定律）]
对任意节点，在任意时刻流入节点的电流代数和为零：
\[
\sum_k I_k=0.
\]
其本质是电荷守恒。
\end{law}

使用节点定律时，关键不是“记忆流入等于流出”，而是先约定每一支路电流正方向，再按符号写出代数和。

\subsection{回路定律}

\begin{law}[基尔霍夫第二定律（回路定律）]
沿任意闭合回路绕行一周，各段电势升降的代数和为零：
\[
\sum_k \Delta V_k=0.
\]
若把电源电动势记为正、电阻上的电压降记为负，则常写为
\[
\sum \varepsilon-\sum IR=0.
\]
其本质是静电场与非静电力共同作用下的能量守恒。
\end{law}

\begin{remark}
回路方程中的符号取决于选定的绕行方向：
\begin{itemize}
  \item 通过电阻时，若绕行方向与电流方向相同，则电势降低，记为 $-IR$；
  \item 通过电源时，若从负极到正极，则电势升高，记为 $+\varepsilon$。
\end{itemize}
只要前后一致，得到的最终物理解不会出错。
\end{remark}

\begin{figure}[H]
\centering
\begin{circuitikz}[scale=1]
  \draw
  (0,0) to[battery1,l=$\varepsilon$] (0,3)
        to[R,l=$R_1$] (3,3)
        to[R,l=$R_2$] (3,0)
        -- (0,0);
\end{circuitikz}
\caption{单回路电路示意图}
\end{figure}

对上图，沿顺时针方向列回路方程得
\[
\varepsilon-IR_1-IR_2=0,
\]
故
\[
I=\frac{\varepsilon}{R_1+R_2}.
\]
这正是串联电阻规律的方程化表达。

例如当图中参数取 $\varepsilon=12\,\text{V}$、$R_1=2\,\Omega$、$R_2=4\,\Omega$ 时，回路定律给出
\[
12-2I-4I=0 \quad\Rightarrow\quad I=2\,\text{A}.
\]
于是两个电阻上的电压降分别为
\[
U_1=IR_1=4\,\text{V},
\qquad
U_2=IR_2=8\,\text{V},
\]
并满足 $U_1+U_2=12\,\text{V}$。

再看一个节点计算：若两条支路分别以 $3\,\text{A}$、$1.5\,\text{A}$ 流入节点，第三条支路流出，则节点方程为
\[
3+1.5-I=0,
\]
故 $I=4.5\,\text{A}$。若第三条支路的参考正方向改设为流入节点，则应写成 $3+1.5+I'=0$，于是 $I'=-4.5\,\text{A}$；负号只表示真实方向与参考方向相反。

\subsection{节点方程与回路方程的联立}

在多支路网络中，未知量往往是若干支路电流。节点定律提供独立的电流关系，回路定律提供独立的电压关系，二者联立即可求解。

\begin{theorem}
对仅含电阻和理想电源的线性电路，节点方程与回路方程组成一组线性代数方程。只要独立方程数与未知支路电流数匹配，便可唯一确定电路状态。
\end{theorem}

这一定理说明：高中电路题的本质通常不是“套技巧”，而是“建立方程”。而微积分的意义在于：当电容、电感加入后，这组方程将从代数方程升级为微分方程。

\section{RC电路}

电容元件会储存电荷，使电流不再瞬间达到稳态。最典型的例子就是 RC 充放电过程。

\subsection{充电方程}

考虑电动势为 $\varepsilon$、电阻为 $R$、电容为 $C$ 的串联电路。设电容器极板上电荷量为 $q(t)$，则电流为
\[
I=\dv{q}{t}.
\]
电容器两端电压为
\[
U_C=\frac{q}{C}.
\]
沿回路应用基尔霍夫第二定律：
\[
\varepsilon-IR-U_C=0.
\]
代入 $I=\dv*{q}{t}$ 得
\[
R\dv{q}{t}+\frac{q}{C}=\varepsilon.
\]
整理为
\[
q'+\frac{1}{RC}q=\frac{\varepsilon}{R}.
\]
这就是 RC 充电的基本微分方程。

\begin{figure}[H]
\centering
\begin{circuitikz}[scale=1]
  \draw
  (0,0) to[battery1,l=$\varepsilon$] (0,3)
        to[R,l=$R$] (3,3)
        to[C,l=$C$] (3,0)
        -- (0,0);
\end{circuitikz}
\caption{RC 串联充电电路}
\end{figure}

\begin{theorem}
初始条件为 $q(0)=0$ 时，RC 充电方程
\[
q'+\frac{1}{RC}q=\frac{\varepsilon}{R}
\]
的解为
\[
q(t)=C\varepsilon\left(1-\mathrm{e}^{-t/(RC)}\right).
\]
因此电流随时间变化为
\[
I(t)=\dv{q}{t}=\frac{\varepsilon}{R}\mathrm{e}^{-t/(RC)}.
\]
\end{theorem}

\textbf{证明思路：} 该方程是一阶线性常微分方程，也可直接验证上述函数满足方程与初值条件。将 $q(t)$ 代入可得
\[
q'(t)=\frac{\varepsilon}{R}\mathrm{e}^{-t/(RC)},
\]
于是
\[
q'(t)+\frac{1}{RC}q(t)
=\frac{\varepsilon}{R}\mathrm{e}^{-t/(RC)}+\frac{1}{RC}\cdot C\varepsilon\left(1-\mathrm{e}^{-t/(RC)}\right)
=\frac{\varepsilon}{R}.
\]
证毕。

\begin{definition}
量
\[
\tau=RC
\]
称为 RC 电路的\emph{时间常数}。当 $t=\tau$ 时，充电电量达到终值 $C\varepsilon$ 的
\[
1-\mathrm{e}^{-1}\approx 63.2\%.
\]
\end{definition}

\begin{remark}
时间常数并不是“充电完成所需时间”，而是衡量系统接近稳态快慢的尺度。通常经历约 $3\tau$ 后已达到终值的 $95\%$ 左右，经历约 $5\tau$ 后可近似视为充电完成。
\end{remark}

\subsection{放电方程}

若电容器初始带电 $q(0)=Q_0$，随后仅通过电阻 $R$ 放电，则回路方程为
\[
IR+\frac{q}{C}=0,
\qquad I=\dv{q}{t}.
\]
故
\[
\dv{q}{t}+\frac{1}{RC}q=0.
\]
其解为
\[
q(t)=Q_0\mathrm{e}^{-t/(RC)},
\qquad
I(t)=-\frac{Q_0}{RC}\mathrm{e}^{-t/(RC)}.
\]
电流中的负号表示放电电流方向与充电时取定的正方向相反。

\begin{figure}[htbp]
\centering
\begin{tikzpicture}
\begin{axis}[
  width=0.75\textwidth, height=0.5\textwidth,
  xlabel={时间 $t/(RC)$}, ylabel={电容电压 $U_C/U_0$},
  axis lines=middle,
  xmin=0, xmax=5.2, ymin=0, ymax=1.08,
  grid=major, grid style={gray!30},
  thick,
  legend style={at={(0.98,0.98)}, anchor=north east, fill=white, draw=gray!30},
]
\addplot[blue, domain=0:5, samples=180] {1 - exp(-x)};
\addplot[red, dashed, domain=0:5, samples=180] {exp(-x)};
\legend{充电, 放电}
\end{axis}
\end{tikzpicture}
\caption{RC 电路中电容两端电压的充电与放电曲线}
\end{figure}

\begin{example}
某 RC 电路中，$R=2\times 10^5\,\Omega$，$C=10\,\mu\text{F}$，电源电动势 $\varepsilon=12\,\text{V}$。求：
\begin{enumerate}
  \item 时间常数；
  \item 充电开始瞬间电流；
  \item $t=2\,\text{s}$ 时电容器上的电荷量。
\end{enumerate}

\textbf{解：}
\[
\tau=RC=2\times 10^5\times 10\times 10^{-6}=2\,\text{s}.
\]
开始瞬间 $t=0$ 时
\[
I(0)=\frac{\varepsilon}{R}=\frac{12}{2\times 10^5}=6.0\times 10^{-5}\,\text{A}.
\]
又因为
\[
q(t)=C\varepsilon\left(1-\mathrm{e}^{-t/\tau}\right),
\]
故当 $t=2\,\text{s}$ 时
\[
q(2)=10\times 10^{-6}\times 12\left(1-\mathrm{e}^{-1}\right)
\approx 7.58\times 10^{-5}\,\text{C}.
\]
\end{example}

\section{RL电路}

电容“阻碍电压突变”，电感则“阻碍电流突变”。在含电感的电路中，电流的变化会产生自感电动势，从而使电流不能突然跳变。

\subsection{电感与自感电动势}

\begin{definition}
若某线圈中电流变化率为 $\dv*{I}{t}$ 时产生自感电动势
\[
\mathcal{E}_L=-L\dv{I}{t},
\]
则称 $L$ 为线圈的\emph{电感}。
\end{definition}

负号体现楞次定律：自感电动势总是阻碍原电流的变化。若电流试图增大，它就产生反向电动势；若电流试图减小，它就“支撑”原电流继续存在。

\subsection{RL 串联电路方程}

考虑电动势为 $\varepsilon$、电阻为 $R$、电感为 $L$ 的串联电路。沿回路列方程：
\[
\varepsilon-L\dv{I}{t}-RI=0.
\]
即
\[
L\dv{I}{t}+RI=\varepsilon.
\]
这就是 RL 电路接通电源后的微分方程。

\begin{figure}[H]
\centering
\begin{circuitikz}[scale=1]
  \draw
  (0,0) to[battery1,l=$\varepsilon$] (0,3)
        to[R,l=$R$] (3,3)
        to[L,l=$L$] (3,0)
        -- (0,0);
\end{circuitikz}
\caption{RL 串联电路}
\end{figure}

\begin{theorem}
若 RL 电路在 $t=0$ 接通电源，且初始电流 $I(0)=0$，则方程
\[
L\dv{I}{t}+RI=\varepsilon
\]
的解为
\[
I(t)=\frac{\varepsilon}{R}\left(1-\mathrm{e}^{-tR/L}\right).
\]
其时间常数为
\[
\tau=\frac{L}{R}.
\]
\end{theorem}

与 RC 充电不同，RC 电路中是\emph{电荷或电压}逐渐升高、电流逐渐减小；而 RL 电路中是\emph{电流}从零逐渐增大到稳态值 $\varepsilon/R$。

\begin{figure}[htbp]
\centering
\begin{tikzpicture}
\begin{axis}[
  width=0.75\textwidth, height=0.5\textwidth,
  xlabel={时间 $t/\tau$}, ylabel={电流 $I/I_{\infty}$},
  axis lines=middle,
  xmin=0, xmax=5.2, ymin=0, ymax=1.08,
  grid=major, grid style={gray!30},
  thick,
]
\addplot[blue, domain=0:5, samples=180] {1 - exp(-x)};
\end{axis}
\end{tikzpicture}
\caption{RL 电路接通后电流按指数规律上升到稳态值}
\end{figure}

若切断电源，仅保留 $R$ 与 $L$ 构成闭合回路，则有
\[
L\dv{I}{t}+RI=0,
\]
解得
\[
I(t)=I_0\mathrm{e}^{-tR/L}.
\]
这说明电感中的磁场能会逐渐释放并在电阻上转化为热。

\begin{example}
某 RL 电路中，$L=0.40\,\text{H}$，$R=20\,\Omega$，$\varepsilon=6\,\text{V}$。求：
\begin{enumerate}
  \item 电路稳态电流；
  \item 时间常数；
  \item $t=0.02\,\text{s}$ 时的电流。
\end{enumerate}

\textbf{解：}
稳态时 $\dv*{I}{t}=0$，故
\[
I_{\infty}=\frac{\varepsilon}{R}=0.30\,\text{A}.
\]
时间常数为
\[
\tau=\frac{L}{R}=0.02\,\text{s}.
\]
于是
\[
I(0.02)=0.30\left(1-\mathrm{e}^{-1}\right)\approx 0.190\,\text{A}.
\]
\end{example}

\begin{remark}
RC 与 RL 两类电路都满足一阶线性微分方程，因此都会呈现指数规律。判断“谁在指数变化”，取决于储能元件：
\begin{itemize}
  \item 电容储存电场能，故 $q$、$U_C$ 常是核心变量；
  \item 电感储存磁场能，故 $I$ 常是核心变量。
\end{itemize}
\end{remark}

更一般地，若把电阻、电感和电容串联起来，回路电流将满足二阶线性方程，并在欠阻尼条件下出现近似形式
\[
I(t)=I_0\mathrm{e}^{-\alpha t}\cos(\omega t+\varphi).
\]
这说明 RLC 电路兼具振荡与耗散两种特征，是机械阻尼振子的电路类比。

\begin{figure}[htbp]
\centering
\begin{tikzpicture}
\begin{axis}[
  width=0.75\textwidth, height=0.5\textwidth,
  xlabel={时间 $t$}, ylabel={电流 $I$},
  axis lines=middle,
  xmin=0, xmax=12, ymin=-1.1, ymax=1.1,
  grid=major, grid style={gray!30},
  thick,
]
\addplot[blue, domain=0:12, samples=220] {exp(-0.22*x)*cos(deg(3.8*x))};
\end{axis}
\end{tikzpicture}
\caption{欠阻尼 RLC 电路中的电流呈指数衰减振荡}
\end{figure}

\section{电功率与焦耳定律}

\subsection{瞬时功率}

单位时间内电场力做的功称为电功率，定义为
\[
P=\dv{W}{t}.
\]
若在时间 $\dd t$ 内有电量 $\dd Q$ 通过某段电路，两端电势差为 $U$，则电场力做功
\[
\dd W=U\,\dd Q.
\]
两边同时除以 $\dd t$ 得
\[
P=U\dv{Q}{t}=UI.
\]

\begin{law}[电功率公式]
电路元件的瞬时功率满足
\[
P=UI.
\]
若元件满足欧姆定律 $U=IR$，则还可写为
\[
P=I^2R=\frac{U^2}{R}.
\]
\end{law}

\subsection{总电能的积分表示}

若电压、电流随时间变化，则在 $[t_1,t_2]$ 时间内电场力做的总功为
\[
W=\int_{t_1}^{t_2} P(t)\,\dd t
=\int_{t_1}^{t_2} U(t)I(t)\,\dd t.
\]
这正是微积分在电路中的一个直接应用：\emph{把瞬时功率沿时间累积，得到总能量。}

\begin{law}[焦耳定律]
纯电阻电路中，电流通过电阻在时间 $t$ 内产生的热量为
\[
Q=I^2Rt.
\]
若电流随时间变化，则应写成积分形式
\[
Q=\int_{t_1}^{t_2} I^2(t)R\,\dd t.
\]
\end{law}

\begin{example}
某电阻两端电压随时间变化为 $U(t)=10\mathrm{e}^{-t}\,\text{V}$，电阻 $R=5\,\Omega$。求从 $t=0$ 到 $t=\infty$ 这段时间内电阻上产生的总热量。

\textbf{解：}
先求电流：
\[
I(t)=\frac{U(t)}{R}=2\mathrm{e}^{-t}\,\text{A}.
\]
故瞬时功率
\[
P(t)=I^2R=4\mathrm{e}^{-2t}\times 5=20\mathrm{e}^{-2t}\,\text{W}.
\]
总热量为
\[
Q=\int_0^{\infty}20\mathrm{e}^{-2t}\,\dd t
=20\cdot \frac{1}{2}=10\,\text{J}.
\]
这里用到了指数函数积分，体现了“瞬时功率积分得到总能量”的思想。
\end{example}

\begin{example}
一只电容器由电源 $\varepsilon$ 通过电阻 $R$ 充电。求整个充电过程中电源提供的总能量，以及最终储存在电容器中的能量。

\textbf{解：}
充电电流为
\[
I(t)=\frac{\varepsilon}{R}\mathrm{e}^{-t/(RC)}.
\]
电源输出功率恒为
\[
P_{\text{源}}=\varepsilon I(t)=\frac{\varepsilon^2}{R}\mathrm{e}^{-t/(RC)}.
\]
故电源提供总能量
\[
W_{\text{源}}=\int_0^{\infty}\frac{\varepsilon^2}{R}\mathrm{e}^{-t/(RC)}\,\dd t
=\frac{\varepsilon^2}{R}\cdot RC=C\varepsilon^2.
\]
最终电容器上电荷量为 $Q=C\varepsilon$，故储能为
\[
W_C=\frac{Q^2}{2C}=\frac{1}{2}C\varepsilon^2.
\]
因此剩余的
\[
W_R=W_{\text{源}}-W_C=\frac{1}{2}C\varepsilon^2
\]
转化为电阻上的焦耳热。这个结论非常经典：RC 充电中，电源提供的一半能量储存在电容器中，另一半在电阻上耗散。
\end{example}

\section{本章小结}

本章建立了稳恒电流与简单暂态电路的基本图景：
\begin{itemize}
  \item 电流是电量对时间的变化率：$I=\dv*{Q}{t}$；
  \item 电流密度的微观表达式为 $\vb{J}=nq\vb{v}_d$；
  \item 欧姆定律既可写成宏观形式 $U=IR$，也可写成微观形式 $\vb{J}=\sigma\vb{E}$；
  \item 基尔霍夫两定律分别体现电荷守恒与能量守恒；
  \item RC、RL 电路都满足一阶线性微分方程，其解都具有指数形式；
  \item 电功率与总能量满足
  \[
  P=UI,\qquad W=\int P\,\dd t,
  \]
  而焦耳热在变电流情形下应采用积分表示。
\end{itemize}

对高中阶段而言，掌握这些公式已经足以解决大量题目；对进一步学习而言，更重要的是看见这些公式背后的结构：\emph{守恒律 + 本构关系 + 微分方程}，这正是理论物理的常见建模方式。

